% framing_success_false_alarm.tex — 2026-08-02
%
% Standalone candidate text describing ONE experiment: on successful
% trajectories only, does key-step-aware or trajectory-propagated
% supervision produce a model more prone to mislabeling a correct
% (safe) prefix as failure? NOT wired into the actual paper build --
% saved here for the author to review/paste in.
%
% Source data: code/keystep_ablation/results/success_false_alarm_curve.csv
% (arm x scope x tau -> false_cut_rate, full threshold sweep),
% success_false_alarm_summary.csv (scalar summaries), and
% success_false_alarm_curves.png (the figure this text references).
% Generated by run_success_false_alarm.py, which loads two versions of
% Ours (macro+micro) -- identical features/folds/calibration/weighting,
% differing only in training label (train_risk_model_temp_sca.py
% --label_mode {keystep,trajectory}) -- and restricts evaluation to
% trajectories that actually succeeded.

\paragraph{Purpose.} Pre-key prefixes of a failed trajectory do not yet
contain the failure symptom, so labeling them as failures -- propagating
a trajectory's terminal outcome to every one of its prefixes -- is a
labeling error, not a defensible alternative target. Rather than argue
this on definitional grounds alone, we test whether training under that
error produces a model measurably more prone to calling a correct
prediction wrong. We isolate this on \emph{successful} trajectories
specifically: their ground truth is identical under both training
schemes (a successful trajectory is labeled 0 at every step regardless
of which label scheme trained the model), so there is no ambiguity about
what the right answer is here -- only about which model gets it right
more often.

\paragraph{Method.} We retrain Ours (macro+micro) with identical
features, folds, calibration, and sample weighting, varying only the
training label: key-step-aware (pre-key prefixes of failed trajectories
relabeled 0) vs.\ trajectory-propagated (every prefix of a failed
trajectory inherits its terminal outcome). Restricting to successful
trajectories from both arms, we sweep a threshold $\tau$ over the full
range of predicted $P(\text{failure})$ and compute, at each $\tau$, the
fraction of successful trajectories that would be flagged at least once
(any step's $P(\text{failure}) > \tau$) -- the same False-Cut-Rate
quantity used elsewhere in this work, reported here across the whole
threshold range rather than at one fixed operating point. We also report
the mean predicted $P(\text{failure})$ across all successful-trajectory
steps as a threshold-free summary.

\begin{figure}[t]
    \centering
    \includegraphics[width=\linewidth]{fig/success_false_alarm_curves.pdf}
    \caption{False Cut Rate on successful trajectories only, swept across
    the full threshold range, by benchmark.}
    \label{fig:success_false_alarm}
\end{figure}

\paragraph{Result.} Figure~\ref{fig:success_false_alarm} shows the
trajectory-propagated model's False-Cut-Rate curve sitting strictly above
the key-step-aware model's at \emph{every} threshold tested, on
\emph{both} WebArena-Lite and Mind2Web individually, with no crossover.
Pooled, mean predicted $P(\text{failure})$ on successful trajectories is
$0.326$ for key-step-aware supervision vs.\ $0.667$ for the
trajectory-propagated ablation -- almost exactly double what it should
be given the ground truth is 0. At a plain $P(\text{failure}) > 0.5$
threshold, $27.6\%$ of successful trajectories are flagged at least once
under key-step-aware supervision, vs.\ $88.2\%$ under the
trajectory-propagated ablation -- nearly nine in ten. The pattern holds
on each benchmark separately (WebArena-Lite: $0.309$ vs.\ $0.636$ mean
$P(\text{failure})$, $19.9\%$ vs.\ $81.0\%$ flagged at $\tau=0.5$;
Mind2Web: $0.336$ vs.\ $0.685$, $32.6\%$ vs.\ $92.9\%$). Propagating the
terminal label to every prefix teaches the model that patterns common to
early success and early failure both mean "failure," and that bias
generalizes directly onto trajectories that never fail at all.
